\begin{portafolioabstract}
	\textbf{Objetivo:} Analizar la diferencia entre redacción libre y redacción académica mediante la reescritura de un texto coloquial sobre el estudio en línea.

	\textbf{Producto elaborado:} Presentación del texto base, identificación de sus rasgos de oralidad e imprecisión y elaboración de una versión académica reformulada en 150--200 palabras.

	\textbf{Enfoque de análisis:} La escritura universitaria no se limita a corregir formas superficiales; exige organizar ideas, jerarquizar argumentos y convertir una experiencia personal en un mensaje claro, interpretable y pertinente para un contexto formativo \citep{nunezcortesEscrituraAcademicaTeoria2015,correaManualBasicoEscritura2018}.
\end{portafolioabstract}

\subsubsection{Introducción}

En redacción académica, la forma de decir una idea es tan importante como la idea misma. Un texto puede comunicar una experiencia válida, pero si está redactado con coloquialismos, saltos lógicos o imprecisión léxica, pierde fuerza argumentativa en contextos universitarios.

Ese es el problema que atiende esta actividad: transformar un mensaje cotidiano sobre el estudio en línea en un texto académico que conserve el significado, pero mejore su estructura, su formalidad y su coherencia. La redacción académica exige claridad, intención comunicativa y articulación de ideas mediante conectores, no solo sustitución de palabras informales \citep{correaManualBasicoEscritura2018,vitoriaEscrituraAcademicaFormacion2019}.

Desde esta perspectiva, el ejercicio de transformación permite observar cómo cambia la eficacia comunicativa cuando el texto pasa de una redacción libre a una redacción orientada por criterios académicos.

\subsubsection{Transformación de redacción libre a redacción académica}

La transformación de redacción libre a redacción académica exige pasar de una expresión espontánea a una construcción discursiva organizada, precisa y orientada por una intención comunicativa explícita. En este marco, escribir no significa solo trasladar ideas, sino jerarquizarlas, vincularlas mediante relaciones lógicas y presentarlas con un registro pertinente para el contexto universitario. Desde esta perspectiva, la diferencia entre ambas formas de escritura no radica en la eliminación de la voz personal, sino en el nivel de control textual que permite convertir una experiencia cotidiana en un argumento claro y comprensible \citep{nunezcortesEscrituraAcademicaTeoria2015,correaManualBasicoEscritura2018,vitoriaEscrituraAcademicaFormacion2019}.

Con base en ello, el ejercicio consiste en transformar un párrafo originalmente redactado en estilo coloquial hacia una versión académica que conserve la idea central, pero mejore su coherencia, su precisión y su fuerza explicativa. El párrafo original es el siguiente:

\portafolioexercisebox{Texto original}{%
``Hoy en día estudiar en línea está medio complicado porque la neta te distraes mucho. O sea, estás leyendo y de repente ya estás viendo redes o contestando mensajes. Además, siento que no entiendes igual que cuando lees en papel, porque como que solo lees rápido y ya. Pero bueno, también es más práctico porque puedes leer en cualquier lado.''%
}

El proceso de transformación en lenguaje académico se desarrolló al identificar la idea central del texto, eliminar marcas de oralidad, reorganizar la secuencia argumentativa y establecer conectores que articularan de manera coherente causas, efectos y valoraciones. También se ajustó el vocabulario para reducir ambigüedades y se revisó la progresión temática con el fin de mantener unidad y claridad en todo el párrafo.

\portafolioexercisebox{Versión académica}{%
El estudio en línea presenta retos vinculados con la gestión de la atención y la calidad de la comprensión lectora. En los entornos digitales, la lectura compite de manera constante con notificaciones, redes sociales y otros estímulos que fragmentan la concentración. Esta interferencia no solo interrumpe el ritmo de trabajo, sino que también puede reducir la profundidad con la que se procesa la información.\par

Además, la lectura en pantalla suele realizarse con mayor rapidez y con menos pausas de reflexión que la lectura en formato impreso, especialmente cuando la persona estudiante no aplica estrategias de autorregulación. Como consecuencia, parte del contenido puede revisarse de manera superficial, lo que afecta la retención y limita el análisis posterior.\par

No obstante, el formato digital ofrece una ventaja operativa clara: permite acceder a materiales desde distintos espacios y dispositivos. Por ello, su eficacia académica no depende únicamente del soporte, sino de las prácticas de lectura y concentración que desarrolle el estudiante. En este sentido, fortalecer hábitos de atención, lectura crítica y manejo del tiempo resulta esencial para aprovechar de manera efectiva la flexibilidad propia de la educación en línea.%
}

Como puede apreciarse en el resultado, la versión académica conserva el contenido esencial del párrafo original, pero mejora su estructura lógica, su formalidad y su capacidad explicativa. La información deja de presentarse como una impresión dispersa y se convierte en una argumentación ordenada, con relaciones explícitas entre problema, consecuencias y valoración final.

\subsubsection{Conclusión}

El ejercicio de transformación permitió comprobar que la escritura académica no se limita a corregir palabras informales, sino a reorganizar el pensamiento para hacerlo más claro, preciso y argumentativamente consistente. La versión original comunicaba una experiencia auténtica, pero su estructura reducía la fuerza del mensaje y dificultaba su proyección en un contexto universitario.

La versión reformulada mantuvo la idea central y mejoró su legibilidad, su coherencia y su capacidad explicativa. En conclusión, transformar redacción libre en redacción académica significa pasar de una expresión espontánea a una construcción discursiva con intención, control y sentido formativo. Ese es, en última instancia, el valor real del ejercicio.
